\section{Introduction}

Riffle shuffling randomizes a deck of cards by cutting it into ordered
piles and combining them by a uniformly random interleaving. In~\cite{SSW},
cutoff was established for general sequences of pile sizes whose empirical
distribution converges to a fixed nondegenerate probability measure on
every macroscopic time interval (of length $\Omega(\log N)$). The mixing
time is then determined by this measure, allowing both deterministic cuts
and random cuts with general distributions.

It is natural to ask what happens when the distribution of the cuts
changes on the mixing time scale. For example, a dealer may first make
nearly balanced cuts and then switch to strongly asymmetric cuts, or use
different distributions during successive intervals of length
$\Theta(\log N)$. Such a process is no longer described by a single
limiting distribution. We study this setting by allowing the distribution
of pile proportions near time $t\log N$ to depend on $t$.

We determine the mixing threshold for these time-inhomogeneous riffle
shuffles under a local nondegeneracy condition. The threshold is described
by an explicit variational expression involving the costs accumulated
over different parts of the shuffling history. As in~\cite{SSW}, there
are two obstructions to mixing: a large class of cards which have followed
the same sequence of piles, and a collection of cold spots where the
cards retain an excess of their initial order. Our result applies to
random pile sizes with arbitrary dependence between different times.

\subsection{Main results}

We first define the shuffling process. Fix $k\ge2$ and consider a deck
of $N$ cards. Given nonnegative integers $\vec n=(n_0,\ldots,n_{k-1})$
with $\sum_i n_i=N$, take the top $n_0$ cards to form the first pile,
the next $n_1$ cards to form the second pile, and so on. We then choose
uniformly among the $\binom{N}{n_0,\ldots,n_{k-1}}$ possible interleavings,
preserving the order within each pile. This defines a Markov transition
matrix $\mathsf P(\vec n)$ on the permutation group $\mathfrak S_N$.
For a deterministic sequence of pile sizes
$\mathbf n=(\vec n^{(j)})_{j\ge1}$, the transition matrix after $K$
shuffles is
\[
 \mathsf P(\mathbf n,K)
 =\mathsf P(\vec n^{(K)})\cdots\mathsf P(\vec n^{(1)}).
\]
Write $P(\mathbf n,K)$ for the resulting distribution when the initial
deck is the identity permutation.

We also consider random sequences of pile sizes
$X=X_N=(X_N^{(j)})_{j\ge1}$, where each $X_N^{(j)}$ is a vector of
nonnegative integers summing to $N$. Conditional on the entire sequence
$X$, the interleavings are chosen independently and uniformly as above.
The distribution of the shuffled deck, averaged over both the pile sizes
and the interleavings, is
\begin{equation}\label{eq:random-shuffle-law}
 P(X,K)=\E_{\mathbf n\sim X}P(\mathbf n,K).
\end{equation}
No independence between the vectors $X_N^{(j)}$ is assumed. Deterministic
sequences are included by taking $X_N$ to be constant. The uniform
distribution on $\mathfrak S_N$, denoted by $\nuUnif$, is stationary for
each shuffle. We write $\dTV$ for total variation distance, and all
logarithms are natural.

To describe the distribution of the cuts, let
\begin{equation}\label{eq:intro-simplex}
 D_k=\left\{p=(p_0,\ldots,p_{k-1})\in[0,1]^k:\sum_i p_i=1\right\},
 \qquad
 V_k=\{p\in D_k:p_i=1\text{ for some }i\}.
\end{equation}
A time profile $\mu_\bullet=(\mu_t)_{t\ge0}$ is a measurable family of
probability measures on $D_k$: for every Borel set $B\subseteq D_k$,
the function $t\mapsto\mu_t(B)$ is measurable. We say that $X$ is
\emph{randomized-$\mu_\bullet$-like} if, for every $T>0$,
\begin{equation}\label{eq:randomized-profile-like}
 \frac1{\log N}\sum_{j=1}^{\lfloor T\log N\rfloor}
       \delta_{(X_N^{(j)}/N,\,j/\log N)}
 \xrightarrow[N\to\infty]{\mathrm{weakly\ in\ probability}}
       \mu_t(dp)\,dt
 \quad\text{on }D_k\times[0,T].
\end{equation}
Here $\delta_x$ denotes a point mass at $x$, and convergence is in the
space of finite measures. Equivalently, for every continuous
$f:D_k\times[0,T]\to\mathbb R$, the corresponding empirical average
converges in probability to $\int_0^T\int f(p,t)\,d\mu_t(p)\,dt$.
When $\mu_t\equiv\mu$, rescaling the time coordinate to $[0,1]$ gives
the randomized-$\mu$-like condition in~\cite{SSW}.

To state the mixing threshold, define for $p\in D_k$ and $\theta\ge1$
\begin{equation}\label{eq:intro-costs}
 a_p=\log\frac1{\max_i p_i},\qquad
 \psi_p(\theta)=-\log\sum_i p_i^\theta,
 \qquad a(t)=\int a_p\,d\mu_t(p).
\end{equation}
We call the profile \emph{locally nondegenerate} if, for every $T>0$,
\begin{equation}\label{eq:profile-nondegeneracy}
 \operatorname*{ess\,inf}_{0\le t\le T}a(t)>0.
\end{equation}
For a constant profile this is exactly the condition $\mu(V_k)<1$.
In general it ensures that, on every bounded time interval, a positive
proportion of the cuts split off two piles of linear size.

For $0\le r\le s$, put
\begin{equation}\label{eq:intro-integrated-costs}
 A(r,s)=\int_r^s a(t)\,dt,\qquad
 B_\theta(r,s)=\int_r^s\int\psi_p(\theta)\,d\mu_t(p)\,dt.
\end{equation}
The two margins are
\begin{align}
 \Lambda_0(T)&=1-A(0,T),\label{eq:intro-largest-margin}\\
 \Lambda_1(T)&=\sup_{0\le s\le u\le T}\inf_{\theta\in[3,4]}
 \left\{\frac{\theta+3}{2}\bigl(1-A(0,s)\bigr)
        -B_\theta(s,u)-2B_2(u,T)\right\}.
 \label{eq:intro-cold-margin}
\end{align}
Here $\Lambda_0$ corresponds to the largest class of equal strings,
while $\Lambda_1$ describes the cold spots. The supremum over $s,u$
allows the latter obstruction to arise during any part of the shuffling
history. Define the critical time
\begin{equation}\label{eq:profile-critical-time}
 C_{\mu_\bullet}
 =\inf\{T>0:\max\{\Lambda_0(T),\Lambda_1(T)\}<0\},
 \qquad \inf\varnothing=+\infty.
\end{equation}

We now state our main result. Randomized-$\mu_\bullet$-like shuffles
have a transition from distance one to distance zero at
$C_{\mu_\bullet}\log N$; in particular, they exhibit cutoff whenever
this critical time is finite.

\begin{theorem}\label{thm:main}
Let $k\ge2$ and let $\mu_\bullet$ be a locally nondegenerate time
profile on $D_k$. For every randomized-$\mu_\bullet$-like sequence
of shuffles $X$ and every fixed $T>0$ with $T\ne C_{\mu_\bullet}$,
\begin{equation}\label{eq:main-profile-limit}
 \lim_{N\to\infty}
 \dTV\bigl(P(X,\lfloor T\log N\rfloor),\nuUnif\bigr)
 =\begin{cases}
   1,&T<C_{\mu_\bullet},\\
   0,&T>C_{\mu_\bullet}.
  \end{cases}
\end{equation}
Consequently, if $C_{\mu_\bullet}<\infty$, then for every
$\varepsilon\in(0,1)$,
\begin{align}
 \lim_{N\to\infty}
 \dTV\bigl(P(X,\lfloor(1-\varepsilon)C_{\mu_\bullet}\log N\rfloor),
                  \nuUnif\bigr)&=1,\label{eq:main-cutoff-lower}\\
 \lim_{N\to\infty}
 \dTV\bigl(P(X,\lfloor(1+\varepsilon)C_{\mu_\bullet}\log N\rfloor),
                  \nuUnif\bigr)&=0.\label{eq:main-cutoff-upper}
\end{align}
If $C_{\mu_\bullet}=+\infty$, the distance tends to one after
$\lfloor T\log N\rfloor$ shuffles for every fixed $T>0$.
\end{theorem}

\subsection{Consequences and Discussion}

For a constant profile $\mu_t\equiv\mu$, the nondegeneracy assumption
reduces to $\mu(V_k)<1$, and Theorem~\ref{thm:main} gives the cutoff
result of~\cite{SSW}. The time-dependent setting also includes
piecewise constant profiles: one may use cuts with empirical
distribution $\mu^{(1)}$ up to time $t_1\log N$, then cuts with
distribution $\mu^{(2)}$ until time $t_2\log N$, and so on. The
functions $A$ and $B_\theta$ retain the costs from each of these
periods in their original order.

The condition~\eqref{eq:randomized-profile-like} only concerns empirical
distributions on macroscopic intervals. It allows deterministic sequences,
as well as random sequences with dependence between different cuts.
For example, within each period of a piecewise constant profile one may
use periodic deterministic cuts or independent random cuts whose
normalized distributions converge to the prescribed measure.

The critical time can be infinite even when the profile is locally
nondegenerate, since the cuts may approach the vertices of $D_k$ as time
increases. A sufficient condition for a finite critical time is
$\int_0^\infty a(t)\,dt=\infty$. Indeed, choosing $\theta=3$ in
\eqref{eq:intro-cold-margin} gives
$\Lambda_1(T)\le3-2A(0,T)$. In particular the critical time is finite
for every constant nondegenerate profile.

The finite estimates underlying Theorem~\ref{thm:main} are uniform in
the prescribed pile sizes. We state them in the upper- and lower-bound
sections, together with profile versions under explicit finite
approximation conditions. The passage from randomized profiles to
these deterministic estimates is given in
Lemma~\ref{lem:randomized-profile-inputs}. The proof of the main theorem
is completed at the end of Section~\ref{sec:lower-bound}.

